% ID: SEL-D0026-OA-E03
% ESTADO: cotejada. Segunda pasada visual completa y borrador compilado revisados.
% Fuente: Selectividad/matematicas_aplicadas/suplemente_septiembre_EXAMEN 4-ANDALUCÍA.pdf
% SHA-256: c41ea8ffe46f2de3a4d7f82e43def97322e29653182ad3df54a57e2ead2beb19
% Página: 1; opción A; ejercicio original 3.

Disponemos de tres monedas: 1 dólar, 1 libra y 1 euro.

La moneda de 1 dólar está trucada y la probabilidad de que salga cara es el doble de la probabilidad de que salga cruz. La moneda de 1 libra también está trucada y tiene dos caras y la de 1 euro es correcta. Se escoge una de las tres monedas al azar y se lanza.
\begin{enumerate}[label=\alph*)]
\item \textbf{(1.5 puntos)} \textquestiondown{}Cuál es la probabilidad de que salga cara?
\item \textbf{(1 punto)} Sabiendo que salió cruz, \textquestiondown{}cuál es la probabilidad de que la moneda lanzada fuera la de 1 dólar?
\end{enumerate}
